\subchapter{Kernel - Installing}{Objective: Install
  the Linux kernel for an ARM target platform.}

\section{Installing kernel using U-Boot}

\subsection{Setting up networking}

This step is to configure U-boot and your workstation to let your
board download files, such as the kernel image and Device Tree Binary
(DTB), using the TFTP protocol through a network connection.

\subsubsection{Ethernet over USB}

Networking will work through the USB device
cable that is already used to power up the board.

{\bf Caution}: For the following to work, make sure that your board
is powered by a USB port on your PC. Otherwise, networking over USB
cannot work.

{\bf Network configuration on the target}

Let's configure networking in U-Boot:

\begin{itemize}
  \item \code{ipaddr}: IP address of the board
  \item \code{serverip}: IP address of the PC host
\end{itemize}

\begin{ubootinput}
=> setenv ipaddr 192.168.0.100
=> setenv serverip 192.168.0.1
\end{ubootinput}

Of course, make sure that this address belongs to a separate network
segment from the one of the main company network.

We also need to configure Ethernet over USB device:
\begin{itemize}
  \item \code{ethprime}: controls which interface gets used first
  \item \code{usbnet_devaddr}: MAC address on the device side
  \item \code{usbnet_hostaddr}: MAC address on the host side
\end{itemize}

\begin{ubootinput}
=> setenv ethprime usb_ether
=> setenv usbnet_devaddr f8:dc:7a:00:00:02
=> setenv usbnet_hostaddr f8:dc:7a:00:00:01
\end{ubootinput}

To make these settings permanent, save the environment:

\begin{ubootinput}
=> saveenv
\end{ubootinput}

% Restart board to reload USB controller or re-bind

% \begin{ubootinput}
% => reset
% \end{ubootinput}

% or

% \begin{ubootinput}
% => unbind ethernet 1
% => bind /ocp/usb@47400000/usb@47401000 usb_ether
% \end{ubootinput}

{\bf Network configuration on the PC host}

To configure your network interface on the workstation side, we need
to know the name of the network interface connected to your board.

Note that when the board is waiting at the U-Boot prompt, no network
interface will show up on the workstation side. It is only when U-Boot
is actively executing a network-related command (such as \code{ping}
or \code{tftp}) that it brings up the USB network connection.

From the board, run \code{ping 192.168.0.1}, and while the \code{ping}
command is running, you should see on your workstation a new network
interface named \code{enx<macaddr>}. Given the value we gave to
\code{usbnet_hostaddr}, it will therefore be
\code{enxf8dc7a000001}. Note that pinging the board from your PC will
not work: when U-Boot is waiting at its prompt, it is not able to
reply to ping requests.

Then, instead of configuring the host IP address from NetworkManager's
graphical interface, let's do it through its command line interface,
which is so much easier to use:

\bashcmd{nmcli con add type ethernet ifname enxf8dc7a000001 ip4 192.168.0.1/24}

\subsubsection{Using ethernet port}
TBD

\subsection{Setting up the TFTP server}

Let's install a TFTP server on your development workstation:

\begin{bashinput}
$ sudo apt install tftpd-hpa
\end{bashinput}

% You can then test the TFTP connection. First, put a small text file in
% the directory exported through TFTP on your development
% workstation. Then, from U-Boot, do:

% \begin{ubootinput}
% => tftp %\zimageboardaddr% textfile.txt
% \end{ubootinput}

% The \code{tftp} command should have downloaded the
% \code{textfile.txt} file from your development workstation into
% the board's memory at location {\tt \zimageboardaddr}\footnote{
% This location is part of the board DRAM. If you want
% to check where this value comes from, you can check the SoC
% datasheet at
% \url{https://www.ti.com/lit/ug/spruh73q/spruh73q.pdf}.
% It's a big document (more than 5,000 pages). In this document, look
% for \code{ARM Cortex-A8 Memory Map} and you will find the SoC memory map.
% You will see that the address range for the memory controller
% ({\em EMIF0 SDRAM})
% starts at the address we are looking for.
% You can also try with other values in the RAM address range.}.

% You can verify that the download was successful by dumping the
% contents of the memory:

% \begin{ubootinput}
% => md %\zimageboardaddr%
% \end{ubootinput}

% We will see in the next labs how to use U-Boot to download, flash and
% boot a kernel.

\subsection{Load and boot the kernel using U-Boot}

% As we are going to boot the Linux kernel from U-Boot,
% we need to set the \code{bootargs} environment corresponding
% to the Linux kernel command line:

% \begin{ubootinput}
% => setenv bootargs console=%\console%
% => saveenv
% \end{ubootinput}
% We will use TFTP to load the kernel image on the board:

% \begin{itemize}

% \item On your workstation, copy the {\tt \ifdefstring{\labboard}{beagleplay}{Image.gz}{zImage}} and DTB
% (\texttt{\dtname}\texttt{.dtb}) to the directory exposed by the TFTP server.

% \item On the target (in the U-Boot prompt), load {\tt \ifdefstring{\labboard}{beagleplay}{Image.gz}{zImage}} from
% TFTP into RAM:
% \begin{ubootinput}
% => tftp %\zimageboardaddr% %\ifdefstring{\labboard}{beagleplay}{Image.gz}{zImage}%
% \end{ubootinput}

% \item Now, also load the DTB file into RAM:
% \begin{ubootinput}
% => tftp %\dtbboardaddr% %\dtname%.dtb
% \end{ubootinput}

% \item Boot the kernel with its device tree:
% \begin{ubootinput}
% => %\ifdefstring{\labboard}{beagleplay}{booti}{bootz} \zimageboardaddr% - %\dtbboardaddr%
% \end{ubootinput}

% \end{itemize}

% \if\defstring{\labboard}{beagleplay}
% This last command should show you an error message of this type:
% \bashcmd{kernel_comp_addr_r or kernel_comp_size is not provided!}

% This is because the boot image that we use, \code{Image.gz}, is compressed, and
% therefore, needs to be uncompressed by U-Boot before continue booting. To do so
% U-Boot needs to know the maximum size of the uncompressed image and where to
% store it.

% If you look at the size of the uncompressed kernel (\code{Image} file),
% you can estimate that 32 MB (0x2000000) is a reasonable upper bound
% for the size of the uncompressed kernel, even with a more exhaustive
% configuration.

% This gives us,

% \begin{ubootinput}
% => setenv kernel_comp_addr_r 0x85000000
% => setenv kernel_comp_size 0x2000000
% => saveenv
% \end{ubootinput}

% Now you can retry the \code{booti} command and see the kernel be uncompressed
% and then loaded.
% \fi

% You should see Linux boot and finally panicking. This is expected: we
% haven't provided a working root filesystem for our device yet.

% You can now automate all this every time the board is booted or
% reset. Reset the board, and customize \code{bootcmd}:

% \begin{ubootinput}
% => setenv bootcmd 'tftp %\zimageboardaddr\ \ifdefstring{\labboard}{beagleplay}{Image.gz}{zImage}; tftp \dtbboardaddr\ {\dtname}.dtb; \ifdefstring{\labboard}{beagleplay}{booti}{bootz} \zimageboardaddr\ - \dtbboardaddr'%
% => saveenv
% \end{ubootinput}

% \ifdefstring{\labboard}{beaglebone}
% {{\bf Known issue}: with at least U-Boot 2023.04 and 2024.04 and with
% USB networking, there is an issue running two \code{tftp} commands in a
% row in \code{bootcmd}. To work around this issue before we get a chance
% to fix it upstream, insert \code{sleep 0.1} between the two \code{tftp}
% commands and everything will work as expected.}{}

% Restart the board to make sure that booting the kernel is now automated.

\section{Installing kernel directly to internal eMMC}
TBD
